\documentclass{article}
\usepackage{enumitem}
\usepackage{amsmath}
\begin{document}

\title{Chapter 33: Biotechnology -- Principles and Processes}
\date{}
\maketitle

\begin{enumerate}[label=Q\arabic*.]

\item Restriction endonucleases cut DNA at specific sequences called:
\\(A) Promoter sequences
\\(B) Palindromic recognition sequences (read same on both strands 5'$\rightarrow$3')
\\(C) Intron sequences
\\(D) Satellite DNA sequences

\item The enzyme reverse transcriptase is used in biotechnology to:
\\(A) Transcribe DNA into RNA
\\(B) Synthesize complementary DNA (cDNA) from an mRNA template
\\(C) Cut DNA at specific restriction sites
\\(D) Ligate DNA fragments

\item In blue-white screening for recombinant clones:
\\(A) Recombinant bacteria appear blue; non-recombinant appear white
\\(B) Non-recombinant bacteria appear blue (lacZ gene intact); recombinant appear white (lacZ disrupted by insert)
\\(C) Both recombinant and non-recombinant bacteria appear white
\\(D) Blue color indicates antibiotic resistance

\item PCR (Polymerase Chain Reaction) requires all of the following EXCEPT:
\\(A) DNA template
\\(B) Primers (short oligonucleotides)
\\(C) Taq DNA polymerase (thermostable)
\\(D) RNA polymerase

\item Assertion: Agarose gel electrophoresis separates DNA fragments by size.\\
Reason: Smaller DNA fragments migrate faster through the gel matrix toward the positive electrode under an electric field.
\\(A) Both true, Reason is correct explanation
\\(B) Both true, Reason is not correct explanation
\\(C) Assertion true, Reason false
\\(D) Assertion false, Reason true

\item A bioreactor differs from a simple fermenter in that:
\\(A) A bioreactor uses microorganisms; a fermenter uses plant cells only
\\(B) A bioreactor provides optimal conditions (pH, temperature, O$_2$, nutrients, mixing) for large-scale production of biologicals
\\(C) A bioreactor operates at very high temperatures
\\(D) A bioreactor does not use microorganisms

\item The Ti plasmid of \textit{Agrobacterium tumefaciens} is used to introduce genes into plants because:
\\(A) It replicates in all plant species independently
\\(B) Its T-DNA integrates stably into the plant nuclear genome during infection
\\(C) It carries antibiotic resistance genes
\\(D) It produces plant hormones

\item Gel electrophoresis staining with ethidium bromide allows visualization of DNA because:
\\(A) EtBr covalently modifies DNA
\\(B) EtBr intercalates between DNA bases and fluoresces under UV light
\\(C) EtBr reacts with the agarose gel
\\(D) EtBr changes the charge of DNA fragments

\item The steps of PCR in correct order are:
\\(A) Annealing $\rightarrow$ Denaturation $\rightarrow$ Extension
\\(B) Denaturation (95°C) $\rightarrow$ Annealing (50--65°C) $\rightarrow$ Extension (72°C)
\\(C) Extension $\rightarrow$ Annealing $\rightarrow$ Denaturation
\\(D) Denaturation $\rightarrow$ Extension $\rightarrow$ Annealing

\item A genomic library differs from a cDNA library in that:
\\(A) A genomic library contains only expressed genes; cDNA library contains all DNA including introns
\\(B) A genomic library contains all DNA sequences including introns and regulatory regions; cDNA library contains only mRNA-derived sequences (no introns)
\\(C) A genomic library uses RNA; cDNA library uses DNA
\\(D) There is no difference between the two

\item Recombinant human insulin produced by bacteria was significant because:
\\(A) It was cheaper than animal insulin only
\\(B) It eliminated the risk of immune reactions to animal insulin and provided an unlimited supply
\\(C) It required no purification
\\(D) Bacterial insulin was more potent than human insulin

\item The selectable marker in a plasmid vector is used to:
\\(A) Provide the origin of replication
\\(B) Identify and select transformed host cells (e.g., antibiotic resistance gene kills non-transformed cells)
\\(C) Allow insertion of foreign DNA
\\(D) Regulate expression of the cloned gene

\item ELISA (Enzyme-Linked Immunosorbent Assay) detects:
\\(A) DNA sequences by hybridization
\\(B) Specific proteins/antigens using antibodies conjugated to enzymes; enzyme catalyzes color reaction indicating positive result
\\(C) RNA by Northern blotting
\\(D) Chromosome abnormalities

\item Match the following biotechnology tools with their uses:\\
1. Restriction endonuclease $\rightarrow$ (a) Cuts DNA at specific sequences\\
2. DNA ligase $\rightarrow$ (b) Joins DNA fragments\\
3. Reverse transcriptase $\rightarrow$ (c) Synthesizes cDNA from mRNA\\
4. Taq polymerase $\rightarrow$ (d) Amplifies DNA in PCR at high temperatures
\\(A) 1-a, 2-b, 3-c, 4-d
\\(B) 1-b, 2-a, 3-d, 4-c
\\(C) 1-c, 2-d, 3-a, 4-b
\\(D) 1-d, 2-c, 3-b, 4-a

\item Downstream processing in biotechnology refers to:
\\(A) The genetic engineering steps to create the recombinant organism
\\(B) Isolation, purification, and formulation of the desired product after biosynthesis
\\(C) Growth of microorganisms in bioreactors
\\(D) Introduction of recombinant DNA into host cells

\end{enumerate}
\end{document}
